\section{评估中的问题与挑战}

\subsection{一致性问题：实现细节的影响}

\begin{figure}[H]
\centering
\includegraphics[width=0.95\textwidth]{slides-images/slide-055.jpg}
\caption{一致性问题：仅将多选题的 ABCD 替换为随机符号，就会改变模型的回答和排名。右侧表格显示 MMLU 的三种不同实现（HELM、Original、Harness）对同一模型给出差异巨大的分数\protect\footnotemark}
\end{figure}
\footnotetext{Slides 第56页。}

MMLU 看似是最简单的评估——四选一多选题。但在近一年的时间里，存在三种主要的 MMLU 实现（HELM、Original、LM Evaluation Harness），它们给出的分数差异巨大：

\begin{itemize}
    \item LLaMA 65B 在 HELM 上为 63.7\%，在 Original 上为 63.6\%，但在 Harness 上仅为 \textbf{48.8\%}
    \item 差异来自两个方面：（1）使用了不同的 prompt；（2）使用了不同的解码策略——约束解码（只在 ABCD 中选）vs 自由生成（可能生成 ``zigzot'' 等无关 token）vs 对数似然比较
\end{itemize}

\begin{warningbox}{``简单''的评估也不简单}
即使是四选一多选题这样看似最简单的评估，实现细节也会导致 15 个百分点的差异。在比较不同论文的基准分数时，必须确认它们使用的是\textbf{同一个评估框架}。这个问题在学界长期存在，很多论文之间的比较可能本就不公平。
\end{warningbox}

\subsection{数据污染（Contamination）}

\begin{figure}[H]
\centering
\includegraphics[width=0.95\textwidth]{slides-images/slide-057.jpg}
\caption{数据污染案例：Horace He 发现 GPT-4 在 2021 年之前的 Codeforces 题目上得分 10/10，但在 2021 年之后的题目上得分 0/10——强烈暗示训练数据包含了测试集\protect\footnotemark}
\end{figure}
\footnotetext{Slides 第58页。}

数据污染是指模型在预训练阶段看过了测试集的数据，本质上是在``测试集上训练''。

\begin{importantbox}{为什么数据污染难以检测}
\begin{enumerate}
    \item 预训练数据量极大，即使有数据访问权限也难以逐一检查
    \item 很多最强的模型是\textbf{闭源}的——你根本不知道它们用了什么训练数据
    \item 效果难以与正常的``能力提升''区分开来
\end{enumerate}
\end{importantbox}

\begin{figure}[H]
\centering
\includegraphics[width=0.95\textwidth]{slides-images/slide-058.jpg}
\caption{标准数据集达到``人类水平''所需的时间在不断缩短——这可能不完全是因为模型真的变好了，也可能是过拟合（overfitting）到固定测试集\protect\footnotemark}
\end{figure}
\footnotetext{Slides 第59页。}

\subsubsection{缓解数据污染的策略}

\begin{figure}[H]
\centering
\includegraphics[width=0.95\textwidth]{slides-images/slide-059.jpg}
\caption{缓解数据污染的方法：私有测试集（如 GSM1K）、动态基准（如 DynaBench）、污染检测技术\protect\footnotemark}
\end{figure}
\footnotetext{Slides 第60页。}

\begin{enumerate}
    \item \textbf{私有测试集}：如 GSM1K——重新采集与 GSM8K 同分布但全新的题目。结果发现开源模型在新数据集上表现显著下降，但 Claude 和 GPT-4 的表现稳定
    \item \textbf{动态基准}（DynaBench）：定期更新测试题目，使预训练数据无法包含未来的测试题。Chatbot Arena 天然具有这个特性
    \item \textbf{污染检测}：
    \begin{itemize}
        \item 观察模型对正确答案的\textbf{置信度}——如果模型对某道题极度确信，可能是在训练中见过
        \item 打乱测试集中\textbf{样本的顺序}——如果模型认为样本 2 应该跟在样本 1 后面（表现为对数似然下降），说明训练数据包含了原始顺序的测试集
    \end{itemize}
\end{enumerate}

\subsection{评估的单一化问题}

\begin{figure}[H]
\centering
\includegraphics[width=0.95\textwidth]{slides-images/slide-061.jpg}
\caption{NLP 基准测试的单一文化：ACL 2021 的 461 篇论文中，70\% 仅评估英语，40\% 仅看准确率\protect\footnotemark}
\end{figure}
\footnotetext{Slides 第62页。}

当前 NLP 评估存在严重的``单一文化''（Monoculture）问题：

\begin{warningbox}{评估的四大盲点}
\begin{enumerate}
    \item \textbf{语言单一}：绝大多数评估仅覆盖英语。BLEU 和 ROUGE 假设可以通过空格分词——在中文、泰文等语言中根本不适用
    \item \textbf{指标单一}：大多数论文仅看准确率，忽略效率、公平性、可解释性等同样重要的维度
    \item \textbf{等权平均}：所有测试样本被赋予相同的权重，忽略了不同样本的价值差异（编写生产代码 vs 推荐餐厅）
    \item \textbf{忽略多样性}：不考虑不同用户群体可能有不同的偏好和需求
\end{enumerate}
\end{warningbox}

\begin{knowledgebox}{多语言基准确实存在}
Mega、Global Bench、Xtreme 等基准涵盖了 30--40 种以上的语言和大量任务。问题不是没有基准，而是学术界\textbf{缺乏使用这些基准的激励}——审稿人通常不会要求评估多语言性能，大家也习惯了只报告英语结果。
\end{knowledgebox}

\subsection{机器翻译中的惯性问题}

\begin{figure}[H]
\centering
\includegraphics[width=0.95\textwidth]{slides-images/slide-063.jpg}
\caption{2019--2020 年机器翻译论文中 82\% 仅报告 BLEU 分数，尽管已知 BLEU 有很多问题且已有更好的替代指标\protect\footnotemark}
\end{figure}
\footnotetext{Slides 第64页。}

\begin{importantbox}{最大的挑战是缺乏改变的激励}
这是``所有挑战中的挑战''——我们知道 BLEU 有问题，我们知道有更好的指标，但 82\% 的机器翻译论文仍然只报告 BLEU。原因是审稿人要求看 BLEU（为了与历史论文可比），这形成了一个自我强化的循环。在学术界，这个问题很难解决；但在实际应用中，如果你知道某个指标不好，\textbf{直接换掉就行}。
\end{importantbox}

\subsection{LLM 评估中的偏见放大}

当使用 GPT-4 等模型进行大规模标注和评估时，其系统性偏差会被放大到整个 NLP 社区：

\begin{figure}[H]
\centering
\includegraphics[width=0.95\textwidth]{slides-images/slide-062.jpg}
\caption{LLM 反映谁的观点：预训练后模型相对均衡（红色），但微调后显著偏向特定群体的观点——这与 RLHF 标注者的人口结构有关\protect\footnotemark}
\end{figure}
\footnotetext{Slides 第63页。}

研究发现，经过 RLHF 微调后的模型倾向于反映标注者群体（通常是东南亚的众包工人和受过高等教育的人群）的偏好和观点。如果整个社区都用同一个有偏差的 LLM 来做评估，这些偏差就会系统性地传播开来。

\begin{itemize}
    \item Anthropic 的 Discover 数据集通过模板化测试了模型在不同种族/性别上的决策偏差（如保险决策）
    \item 结果表明，某些群体确实受到了更多的歧视——尽管这并不令人意外
    \item MLPerf 等基准开始关注计算效率——在给定时间内达到目标性能
\end{itemize}

\subsection{本章小结}

当前 LLM 评估面临五大挑战：（1）实现细节导致的一致性问题——即使是 MMLU 这样``简单''的基准，不同实现也可能给出相差 15\% 的分数；（2）数据污染——闭源模型的训练数据不透明；（3）评估的语言和指标单一化；（4）学术界缺乏改变评估方式的激励；（5）LLM 评估可能放大系统性偏见。

%% ========== 总结与延伸 ==========

